% Source: Wald GR, physical PDF p. 332
%         (printed p. 319).
% Coverage: Side and top views of a Kerr ergosphere.
% Figures/tables/footnotes: Figure 12.6; no tables or footnotes.
% Status: source-reviewed and compile-clean.
% Uncertainties: none.

\begingroup
\renewcommand{\thefigure}{12.6}
\begin{figure}[H]
  \centering
  \begin{tikzpicture}[x=.75cm,y=.75cm,>=Latex,line cap=round,line join=round,font=\small]
    % Side view.
    \draw[thick] (-4.0,0) ellipse [x radius=1.72,y radius=.67];
    \draw[thick] (-4.0,0) circle [radius=.72];
    \draw[thick] (-4.0,.20) -- (-4.0,1.70);
    \draw[->,thick] (-4.28,1.45) arc[start angle=200,end angle=520,x radius=.28,y radius=.11];
    \draw[->] (-2.58,.18) .. controls (-3.05,.16) and (-3.15,.06) .. (-3.34,.00);
    \node[right] at (-2.47,.44) {Black Hole};
    \draw[->] (-2.55,-1.05) .. controls (-3.05,-.82) and (-3.22,-.50) .. (-3.30,-.33);
    \node[right] at (-2.42,-1.08) {Ergosphere};
    \node at (-4.0,-1.85) {(a)};
    % Top view.
    \draw[thick] (3.0,0) circle [radius=1.37];
    \draw[thick] (3.0,0) circle [radius=.72];
    \draw[->,thick] (2.55,.15) arc[start angle=170,end angle=35,radius=.46];
    \node at (3.0,-1.85) {(b)};
  \end{tikzpicture}
  \caption{A sketch showing (a) a \lq\lq side view\rq\rq\ and (b) a \lq\lq top view\rq\rq\ of
  the ergosphere of a Kerr black hole.}
  \label{fig:12.6}
\end{figure}
\endgroup
